% a01_H0rel part 5: the perturbative input and the lower bound
\subsection{The perturbative input on the cube}\label{app:a01:TD}

Ba{\l}aban's perturbative bounds \cite[Prop.~1]{Ba3} are consumed, in the lower bound at every step and in the upper bound at every scale, in the field-dependent form: for a renormalised family $\mathcal R$ of order $\le\bar n=13$ of the expansion of the effective action, built from the lines (slices $G_{(i)}$ and $G_k$) and legs ($\phi^{(k)}$, photon backgrounds) of the step, with external fields in the class of \cref{lem:a01:legs} and possibly source vertices, the sum over orderings, indices and positions of the values of its members, with one vertex $x_0$ held, is bounded by $C\ell_k^{\beta_0}p(\ell_k)^{n_1}(1+K-k)^{p_0}e^{-\delta\dist(V_*)}$ times the external-field norms, with $V_*$ the set of vertices carrying external legs or sources and $\beta_0=\frac14-\beta$; the $R$-terms (orders $>\bar n$) are $O(1)\ell_k^{\kappa'}|\Lambda^{(k)}|$ with $\kappa'>d$. On the torus this is \cite[Prop.~1]{Ba3} with the hypotheses on the external fields supplied by \cite[Lemma~2.4]{Ba2}.

\begin{proposition}[the perturbative bounds on the cube]\label{prop:a01:TD}
On $\CK$ with King's flow, at every scale $k\le K_f$ and every regular $A$: (a) on the whole cube, the bound above holds for every renormalised family with the cube's lines and legs, with the counterterm King's free-boundary $\dm$ (\cite[p.~330]{K2}, field- and domain-independent), in the organisation of \cite[p.~330]{K2}: for a held vertex $x_0$ at face distance $t(x_0)\ge1$ (zone I) with the torus constants; for $t(x_0)<1$ (zone B) with the additional factor $\Lambda(x_0)\le C(1+\ln_+(1/t(x_0)))^{14}t(x_0)^{-(1-\beta)}e^{-\delta_\Lambda t(x_0)}$, which is summable: $\sum_{x_0\in\text{unit cube}}\vth^d\Lambda(x_0)\le C$; (b) the same on every sub-domain $\Omega$ of the cube (any union of $k$-blocks, meeting $\DD$ or not) for the families with the lines $G_{(i)}(\Omega,A)$, $G_k(\Omega,A)$ and the legs $\phi^{(k)}_\Omega$, and with one line replaced by a difference line $G_k(\Omega,A)-G_k(\Omega_2,A)$ of two Neumann domains as an external pair, provided every vertex cube lies at depth $\ge R_*$ in every propagator domain of the family; in the difference case the bound carries the extra factor $\exp[-\delta_c(\dist(\mathrm{cube}(v),\Xl\setminus\Omega_2)+\dist(\mathrm{cube}(v'),\Xl\setminus\Omega_2))]$ for the two ends $v,v'$ of the difference line; (c) the consequences \cite[(3.58)]{Ba1} hold with these inputs: the coefficients of the effective potential $V^{(k)}$ are $O(1)\ell_k^{\beta_0}$ and the remainders $O(1)\ell_k^{\kappa'}|\Lambda^{(k)}|$, uniformly in $k$, $K$, $\vth$ and the field.
\end{proposition}
\begin{proof}
(a) Zone I is the torus bound: at face distance $\ge1$ every line is, up to the face-localised difference of \cref{lem:a01:FLD} (exponentially small in the distance to the face and absorbed in the decay factors), the torus line at the same field, and the photon lines are the image sums of \cref{lem:a01:NL}; the class of the external legs is \cref{lem:a01:legs}. Zone B is the content of the boundary-remainder bounds established in the proof of \cref{prop:k01:Kpos} (the scalar boundary remainders at zero background and the photon face terms of \cite[pp.~331--332]{K2}, with the slice bounds $|D^{n_x}_xD^{n_y}_yG^\Omega_{(i)}(b,b')|\le C\ell_i^{-1-n_x-n_y}e^{-\kappa|x-y|/(2\ell_i)}$ and the cumulative line bound $|D^{n_x}_xD^{n_y}_yG^\Omega_{\le i_0}(b,b')|\le C|x-y|_\vth^{-1-n_x-n_y}e^{-\kappa|x-y|/(4\ell_{i_0})}$ at zero background on the cube and on King's localisation boxes, face bonds included, which follow from the odd-image representation of Step~1 of \cref{lem:a01:FL} and the free weighted bounds; and, at regular $A$, from \cref{lem:a01:A23,lem:a01:A4,lem:a01:A5} through the gauge-away of the constant part on each box): there every line is written as its free counterpart plus a boundary remainder localised at the face, and the remainder vertices are of higher degree, so that the degree bookkeeping of \cite[Sect.~2]{Ba3} is unchanged while the sum over the position of the held vertex produces $\Lambda(x_0)$. The photon-image part of the counterterm is the photon-image part of the same remainder families (\cref{app:a01:photon}). (b) is \cref{thm:a01:B2} below. (c) is \cite[p.~421]{Ba3}: the coefficient bounds use of the external scalar field exactly the sup clause, the gradient clause (P-grad) and the block clause (P-block) of \cref{lem:a01:PC}, at the current vertex \cite[(1.8)]{Ba3} and at the first two kernel vertices \cite[(1.18)]{Ba3}, every other family using the sup clause alone; on the cube the current vertex sits on live bonds only (a face bond carries no link) and the kernel vertices at weighted sites only (the weight vanishes at killed sites), so the two clauses are needed only there.
\end{proof}

